\documentclass[a4paper,12pt]{article}
\usepackage[a4paper,margin=1in]{geometry}
\usepackage{enumitem}
\usepackage{hyperref}
\usepackage[official]{eurosym}
\usepackage{textcomp}
\newenvironment{benumerate}[1]{
    \let\oldItem\item
    \def\item{\addtocounter{enumi}{-2}\oldItem}
    \begin{enumerate}[leftmargin=*,noitemsep]
    \setcounter{enumi}{#1}
    \addtocounter{enumi}{1}
}{
    \end{enumerate}
}

\setlength{\parindent}{0pt}
\setlength{\parskip}{0.5em}

\usepackage[T1]{fontenc}
\usepackage{lmodern}
\usepackage{titlesec}

% A simple CV header
\newcommand{\cvheader}[5]{%
  {\LARGE\bfseries #1}\par
  \vspace{0.25em}
  {\small #2 \;|\; #3 \;|\; \href{#4}{#4} \;|\; #5}\par
  \vspace{0.75em}
  \hrule
  \vspace{1.0em}
}

\begin{document}


\section{Personal Information}
% Full name, Birth date (YYYY-MM-DD), Home address, Phone number, Email
Full name:  Han Wang\\
Work Address: Bundesstr. 53, 20146, Hamburg, Germany\\
Phone Number: +49 40 42838-5384 \\
Email: han.wang@uni-hamburg.de or hannnwangus@gmail.com\\
Personal Website: \url{https://hannnwang.github.io/}\\
ORCID: \url{https://orcid.org/0000-0002-5841-5474}\\
Language: Mandarin Chinese (native), English (fluent), German (beginner)

\section{Degrees}
\begin{itemize}[leftmargin=*,noitemsep]
\item 08/2011-07/2015  B.S. in Atmospheric Sciences, University of Science and Technology of China,
China. Dissertation advisor: Prof. Dr. Rui Li
\item 09/2015 -07/2020   PhD in Atmosphere-Ocean Science and Mathematics, Courant Institute of Mathematical Sciences, New York University, United States. Dissertation advisor: Prof. Dr. Oliver B{\"u}hler
\end{itemize}
\section{Employment History }
\begin{itemize}[leftmargin=*,noitemsep]
\item 09/2024-present  Project Leader, Theoretical Oceanography, University of Hamburg
\item 04/2022-09/2024  Postdoctoral Research Associate, School of Mathematics, University of Edinburgh
\item 09/2020-03/2022   Postdoctoral Researcher, Department of Physics, University of Toronto
\end{itemize}
\textbf{Leaves}
\begin{itemize}[leftmargin=*,noitemsep]
\item 07/2020 – 09/2020: Gap due to visa delays caused by Canadian governmental closures during COVID-19
\end{itemize}

\section{Research Interests}
\textbf{Physical Oceanography and Fluid Dynamics}
\begin{itemize}[leftmargin=*,noitemsep]
\item{Wave-mean interactions and disentanglement}
\item{Submesoscale ocean dynamics}
%\item[--]{Analysis of observational data}
\item{Statistical fluid mechanics}
\item{Deep learning}
\end{itemize}

\section{Scientific Expertise}
\subsection{Peer-reviewed Journal Articles}
\begin{itemize}
\item \textbf{Wang, H.}, B\^{o}as, A.B.V.,  Vanneste J. and Young, W.R. \textit{The U2H map explains the effect of (sub)mesoscale turbulence on significant wave height statistics}. In press (2025) on Journal of Physical Oceanography. \href{https://arxiv.org/abs/2504.21736}{[access]}
\item \textbf{Wang, H.}, B\^{o}as, A.B.V.,  Vanneste J. and Young, W.R. (2025). \textit{Scattering of surface waves by ocean currents: the U2H map}. Journal of Fluid Mechanics,  1005. \href{https://arxiv.org/abs/2402.05652}{[access]}
%\href{https://www.youtube.com/watch?v=migsRGIAB-c}{[talk]}
\item \textbf{Wang, H.}, B\^{o}as, A.B.V., Young, W.R. and Vanneste J. (2023). \textit{Scattering of swell by currents}. Journal of Fluid Mechanics, 975.  \href{https://arxiv.org/abs/2305.12163}{[access]}
\item \textbf{Wang, H.}, Grisouard, N., Salehipour, H., Nuz, A., Poon, M., and Ponte, A. L. (2022). \textit{A deep learning approach to extract internal tides scattered by geostrophic turbulence}. Geophysical Research Letters, 49(11), e2022GL099400. \href{https://essopenarchive.org/doi/full/10.1002/essoar.10508849.2}{[access]}
\item Khatri, H., Griffies, S. M., Uchida, T., \textbf{Wang, H.},  and Menemenlis, D. (2021). \textit{Role of mixed-layer instabilities in the seasonal evolution of eddy kinetic energy spectra in a global submesoscale permitting simulation.} Geophysical Research Letters, 48(18), e2021GL094777. \href{https://repository.library.noaa.gov/view/noaa/43143}{[access]}
\item \textbf{Wang, H.} and B{\"u}hler, O. (2021). \textit{Anisotropic statistics of Lagrangian structure functions and Helmholtz decomposition.} Journal of Physical Oceanography, 51(5), 1375-1393. \href{https://hannnwang.github.io/WangBuhlerJPO21.pdf}{[access]}
\item \textbf{Wang, H.} and B{\"u}hler, O. (2020). \textit{Ageostrophic corrections for power spectra and wave–vortex decomposition.} Journal of Fluid Mechanics, 882. \href{https://hannnwang.github.io/AgeoSpectraJFM20.pdf}{[access]}
\item[] -- Highlighted in \href{https://www.cambridge.org/core/journals/journal-of-fluid-mechanics/article/untangling-waves-and-vortices-in-the-atmospheric-kinetic-energy-spectra/1BEB1ABC32CD2CFAB99BAFEE4712CD0C}{Focus on Fluids}
\end{itemize}
\subsection{Submitted Manuscripts}
\begin{itemize}
\item \textbf{Wang, H.}, Uncu J., Srinivasan K., Grisouard, N. \textit{Disentangling internal tides from balanced motions with deep learning and surface field synergy}. In Revision Journal of Advances in Modeling Earth Systems. \href{https://arxiv.org/html/2511.03614v1}{[access]}
\end{itemize}

\subsection{Research Grants}
Received as principal applicant:
\begin{itemize}[leftmargin=*,noitemsep]
\item Project W2: ``Observed and simulated internal tides: generation, modification by eddies, and contribution to energy budget'' of the Collaborative Research Centre TRR 181:  ``Energy transfers in Atmosphere and Ocean''. Deutsche Forschungsgemeinschaft (Projektnummer 274762653).  Award date  06/2024. Funded period 07/2024 - 06/2028. Funded value \euro 621,100. 
\end{itemize}
Named as intended postdoctoral researcher (not PI):
\begin{itemize}[leftmargin=*,noitemsep]
\item ``Disentangling internal waves and (sub–)mesoscale motions in satellite altimetry: Northeast Pacific”. Canadian Space Agency.  Award date 12/2023.  Funded period 04/2024 – 03/2027 (estimated). Funded value \$225,000. The other applicants are: Jody Klymak (principal applicant), Guoqi Han (co-applicant), Tetjana Ross (co-applicant). 
\item``Phase Averaged Deferred Correction for Multi-Timescale Systems”. UK Research and Innovation (Project Reference: EP/Y032624/1). Funded period 06/2024 – 02/2025. Funded value \textsterling78,966. The other applicant is Hossein Kafiabad (principal applicant). 
\end{itemize}

\section{Presentations}
\subsection{Conference Presentations}
\textbf{Oral Presentations}
\begin{itemize}[leftmargin=*,noitemsep]
\item Born\"{o} workshop on surface waves,  Stora Born\"{o}, Sweden,  `` The  U2H map  explains effects of (sub)mesoscale currents on significant wave height'', Aug 2025
\item GFD Symposium in memory of Vladimir Zeitlin, Paris, France, ``Synergizing Surface Fields in a Deep-learning Extraction of Internal Tides'', May 2025.
\item TAPGFD (Theoretical and Practical Perspectives in Geophysical Fluid Dynamics, Bengaluru, India, virtual presentation, ``A deep learning approach to extract surface internal tidal signals scattered by geostrophic turbulence”, May 2024.
\item Ocean Sciences Meeting, New Orleans,  United States, ``Imprint of Currents on Surface Waves: solutions in two regimes ”, Feb 2024.
\item WISE Zoominar on Waves over spatial inhomogeneity, virtual.``Imprint of Currents on Surface Waves", Oct 2023.
\item  Challenger Society Ocean Modelling Conference, Southampton, United Kingdom, ``Dynamical insights from Lagrangian-filtered structure functions applied to drifter observations”, Sep 2023.
\item EGU General Assembly, Vienna, Austria, ``Imprint of ocean currents on significant wave height”, Apr 2023.
\item  TRR 181 Eddy-Wave Meeting, Hamburg, Germany, virtual presentation, ``Dynamical insights from frequency-filtered Lagrangian structure functions", Feb 2023.
\item 103rd American Meteorological Society Annual Meeting (23rd Conference on Air-Sea Interaction), United States, virtual presentation, ``Imprint of ocean currents on significant wave height", Jan 2023.
\item Oberwolfach Workshop 2238 - Multiscale Wave-Turbulence Dynamics in the Atmosphere and Ocean, Oberwolfach, Germany, ``A deep learning approach to extract surface internal tidal signals scattered by geostrophic turbulence”, Sep 2022.
\item IX International Symposium on Stratified Flows, Cambridge, UK,``A deep learning approach to extract surface internal tidal signals scattered by geostrophic turbulence”, Aug 2022.
\item Surface Water and Ocean Topography (SWOT) Science Team Meeting, virtual presentation,``Internal tidal extraction: challenges from scatterings by vortices, and hopes for a deep learning solution”, Jun 2022.
\item Ocean Sciences Meeting, virtual, ``Extraction of tidal signals from a machine learning approach”, Apr 2022.
\item EGU General Assembly, virtual, ``Anisotropic statistics of Lagrangian structure functions and Helmholtz decomposition”, Apr 2021.
\item Meeting on eddies and internal waves with TRR Mercator fellows, TRR 181, virtual, ``Generalizing the “BCF14” method to anisotropic cases”, Mar 2021.
\end{itemize}
\textbf{Poster Presentations}
\begin{itemize}[leftmargin=*,noitemsep]
\item 12th Warnem\"{u}nde Turbulence Days (WTD) on ``Waves and Turbulence", Insel Vilm, Germany, ``Synergizing Surface Fields in a Deep-learning Extraction of Internal Tides'', Sep 2025.
\item  Climate Exploration in Lively Liaison with the Ocean (CELLO), Hamburg, Germany, ``The  U2H map  explains effects of (sub)mesoscale currents on significant wave height'', Sep 2025
\item 22nd Conference on Atmospheric and Oceanic Fluid Dynamics, Maine, UNITED STATES, ``Anisotropic Helmholtz decomposition of Lagrangian Tracer Data”, Jun 2019.
\item 21st Conference on Atmospheric and Oceanic Fluid Dynamics, Oregon, United States, ``Wave-Vortex Decomposition of 1D ship-track data with weak nonlinearity in the balanced flow”, Jun 2017.
\end{itemize}
\subsection{Invited Seminars and Research Visits }
\begin{itemize}[leftmargin=*,noitemsep]
\item University of Bremen, Bremen, Germany, ``Diagnosing scale-dependent dynamics from limited observations”. Jan 2025.
\item University of Edinburgh, Edinburgh, United Kingdom. Research visit only (no talk given). Nov 2024.
\item TIFR Centre for Applicable Mathematics, Bengaluru, India, virtual presentation, ``A deep learning approach to extract internal tides scattered by geostrophic turbulence”. Oct 2024.
\item Durham University, United Kingdom, ``Scattering of surface waves by oceanic currents”. Jun 2024.
\item University of California, Los Angeles, California, United States. Research visit only (no talk given). Mar 2024.
\item California Institute of Technology , California, United States, ``Scattering of surface waves by oceanic currents”. Mar 2024.
\item University of California San Diego, California, United States, ``Scattering of surface waves by oceanic currents”. Mar 2024.
\item Colorado School of Mines, Colorado, United States, ``Scattering of surface waves by oceanic currents”. Feb 2024.
\item University of Hamburg, Hamburg, Germany,``A deep learning approach to extract internal tides scattered by geostrophic turbulence”. Jul 2023.
\item University of Waterloo, Canada, ``Disentangling balanced and unbalanced flows under weak nonlinearity”, virtual,  Sep 2021.
\end{itemize}


\end{document}